%!TEX root = ../../main.tex
%% ===============================================================
%% 8.6 재현성
%% 파일: chapters/08_experimental_setup/06_reproducibility.tex
%% 상위: chapters/08_experimental_setup/chapter.tex
%% 정의 라벨: sec:setup-repro
%% 참조 라벨: -
%% 포함 객체: -
%% 인용: -
%% 상태: 초안 (docs/OUTLINE.md 참고)
%% ===============================================================

\section{재현성}
\label{sec:setup-repro}

난수 시드는 NumPy·PyTorch·CUDA에 걸쳐 고정하며, 분할·캐시는 버전 태그로 무효화된다. 모든 설정은 단일 데이터클래스에 모여 있고 실험 결과와 함께 직렬화된다. 파이프라인의 계약(무누출 보간, 국소화 상수, 특징 차원, 라벨 가중치)은 \numTests{} 개의 단위 시험으로 고정되어 있다. 논문에 대응하는 코드 상태는 태그 \texttt{v1.0-cog2026}으로 공개되어 있으며, 완전한 실험은 \texttt{run\_paper\_full.py}(3 시드 $\times$ 2 단계)로 재실행할 수 있다. GPU 연산의 일부(예: 산란 연산)는 결정론이 강제되지 않으므로 시드 간 편차가 신뢰구간에 반영된다.
